% uid: ""
% school: "lycee"
% level: "3eme"
% section: ["Economie et Gestion"]
% chapter: "Statistiques"
% title: "Un salaire extrême : moyenne ou médiane ?"
% difficulty: "Moyen"
% points: null
% tags: ["statistiques", "robustesse", "valeur-extreme", "coefficient-de-variation", "mediane"]
% provenance: "own"
% dressing: "valeur extrême"
% rung: 3
% batch: "2026-11-20-stat-3eco-A"

% ---STATEMENT---
Dans une petite entreprise, les $10$ salaires mensuels (en dinars) sont :
$$420\ ;\ 450\ ;\ 450\ ;\ 480\ ;\ 500\ ;\ 520\ ;\ 540\ ;\ 560\ ;\ 600\ ;\ 2\,000.$$
Le dernier est celui du directeur. On note $S_1$ la série complète et $S_2$ la série des $9$ autres salaires.

\begin{enumerate}
  \item Pour $S_1$ : calculer la moyenne, la médiane, les quartiles, l'écart type et le coefficient de variation $CV=\dfrac{\sigma}{\overline{x}}$.
  \item Mêmes calculs pour $S_2$.
  \item Compléter : le salaire du directeur fait passer la moyenne de \ldots à \ldots, alors que la médiane ne varie que de \ldots à \ldots.
  \item Pour présenter « le salaire typique » à un candidat, quel indicateur de position choisir pour $S_1$ ? Justifier.
\end{enumerate}

% ---SOLUTION---
\begin{enumerate}
  \item \textbf{Série complète $S_1$} ($N=10$).

  \emph{Moyenne} : $\sum x_i=6\,520$, donc $\overline{x}=652$.

  \emph{Médiane} : rangs $5$ et $6$ : $500$ et $520$, donc $M_e=510$.

  \emph{Quartiles} : $\dfrac{10}4=2{,}5\to$ rang $3$ : $Q_1=450$ ; $\dfrac{3\times10}4=7{,}5\to$ rang $8$ : $Q_3=560$.

  \emph{Écart type} : $\sum x_i^2=6\,297\,400$ (le seul carré $2\,000^2=4\,000\,000$ en pèse plus de la moitié), donc
  $$V=\frac{6\,297\,400}{10}-652^2=629\,740-425\,104=204\,636\qquad\sigma\approx452{,}4.$$
  \emph{Coefficient de variation} : $CV=\dfrac{452{,}4}{652}\approx$ \ib{$0{,}69$}.
  \item \textbf{Série $S_2$} ($N=9$ : on retire $2\,000$).

  \emph{Moyenne} : $\dfrac{4\,520}9\approx502{,}2$. \emph{Médiane} : rang $5$, donc $M_e=500$.

  \emph{Quartiles} : $\dfrac94=2{,}25\to$ rang $3$ : $Q_1=450$ ; $\dfrac{27}4=6{,}75\to$ rang $7$ : $Q_3=540$.

  \emph{Écart type} : $\sum x_i^2=2\,297\,400$ donc $V=\dfrac{2\,297\,400}{9}-\left(\dfrac{4\,520}9\right)^2\approx255\,266{,}7-252\,227{,}2\approx3\,039{,}5$ et $\sigma\approx55{,}1$.

  \ib{$CV\approx\dfrac{55{,}1}{502{,}2}\approx0{,}11$}.
  \item Un seul salaire fait passer la moyenne de $502{,}2$ à $652$ (hausse de près de $150$ dinars), la médiane de $500$ à $510$ seulement. L'écart type passe de $55$ à $452$ : il est multiplié par plus de $8$.
  \item \textbf{Rappel.}\par
  \begin{prop}[Choisir un indicateur de position]
  La \emph{moyenne} utilise toutes les valeurs : une valeur extrême la déplace beaucoup. La \emph{médiane} ne dépend que du rang : elle est \emph{robuste}. En présence d'une valeur extrême, on préfère le couple $(M_e\,;Q_3-Q_1)$ à $(\overline{x}\,;\sigma)$.
  \end{prop}
  Pour présenter « le salaire typique » on choisit donc la \textbf{médiane} ($510$ dinars) : $9$ salaires sur $10$ sont inférieurs ou égaux à $600$. La moyenne de $652$ dépasse $9$ salaires sur $10$ : elle ne décrit personne.

  \textbf{Piège :} cela ne veut pas dire que la moyenne est fausse, mais qu'elle répond à une autre question (« combien l'entreprise verse-t-elle en moyenne par salarié ? »).
\end{enumerate}

\textbf{Idée :} une valeur extrême tire la moyenne et l'écart type, mais presque pas la médiane : on choisit l'indicateur selon la série.
